\section{Area positivity and preparation of the asymptotic end}
\label{prep:section}
\label{sec:preparation}

This section makes the reductions needed before the scalar construction.
The reductions preserve the electric flux form and the prescribed trapping
sign on each original boundary component.  They do not require a filling,
a preferred slice, or a timelike ADM vector at the outset.  The conclusion
about the original ADM vector will follow only after the estimate for
prepared data has been proved in Proposition~\ref{flux:prepared-bound}.
Throughout this section, $k=n-1$, $\omega=\Area(S^k)$, and
$c_n^2=k(k-1)/2$.  Set
\[
 \mu=\frac12\bigl(R_g+(\tr_gK)^2-|K|_g^2\bigr),
 \qquad J=\operatorname{div}_g(K-(\tr_gK)g),
 \qquad \alpha=i_E\vol_g.
\]
Thus $\dd\alpha=0$ on the original exterior.  A smooth positive radius
function, denoted by $r$, is fixed on the entire exterior and agrees with
the coordinate radius outside a compact set.  Constants on the compact
part may depend on this extension.

\subsection{The full-cut convention and positive area}

\begin{lemma}[A bounded flux form]
\label{prep:positive-area}
Under the topological, completeness, and asymptotic hypotheses of
Definition~\ref{def:data}, the full-cut infimum satisfies
\[
 0<a\leq\Area_g(S).
\]
This assertion does not use the constraint inequalities or an assumed
minimizing hypersurface.
\end{lemma}

\begin{proof}
Choose a small boundary disk in one component of $S$.  A smooth embedded
arc beginning transversely at that disk can be continued into the
coordinate end and then along a straight coordinate ray.  By shrinking a
tubular neighborhood, obtain a proper embedded tube
\[
 [0,\infty)\times B^{n-1}_\delta\longrightarrow\Omega
\]
whose initial face lies in the chosen boundary disk.  The transverse
coordinates have fixed small coordinate width sufficiently far along the
ray.  Choose a nonnegative function
$\zeta\in C_c^\infty(B^{n-1}_\delta)$ with integral one.  The form
\[
 \zeta(z)\,\dd z^1\wedge\cdots\wedge\dd z^{n-1}
\]
on the tube is closed.  Extend it by zero across the lateral boundary;
call the resulting smooth closed form $\beta$.  Its comass is bounded:
on the compact part this follows from smoothness, and on the end from
asymptotic Euclidean comparability and the fixed transverse width.
Orient the initial disk so that $\int_S\beta=1$.

Every full cut $\Gamma$ is homologous to $S$ with the orientation toward
the end.  This can be seen by truncating both exteriors at one distant
coordinate sphere and applying Stokes to their difference; portions of
$S$ shared by the two boundaries cancel with their full multiplicity.
Equivalently, integration over the compact difference of the two
codimension-zero sets gives the same identity.  Thus
\[
 1=\int_\Gamma\beta
 \leq \|\beta\|_{\mathrm{comass},\infty}\Area_g(\Gamma).
\]
Taking the infimum proves the lower bound.  The admissible choice
$\Gamma=S$ gives the upper bound.  In particular, this argument uses the
entire intrinsic boundary and never replaces it by an empty relative
frontier.
\end{proof}

The same Stokes comparison, applied to $\alpha$, shows that every full
cut has signed flux $\omega Q$.  No sign condition is imposed on the
flux of an individual component.  Fixing $\alpha$ in all the following
constructions therefore fixes both the signed total charge and all
original boundary-component fluxes.

\subsection{Conformal changes and a conditional endpoint reduction}

We will also need a limited auxiliary magnetic class.  Let $D_m$ be a
smooth closed two-form of compact support, normalized by dividing the
spatial Faraday form by $c_n$.  Keep both $\alpha$ and $D_m$ fixed when
changing the metric.  For $|v|_g\leq1$, define
\begin{equation}
 \label{prep:electromagnetic-cost}
 I_g(v)=c_n^2\bigl(|E|_g^2+|D_m|_g^2+2D_m(v,E)\bigr),
 \qquad \mu_m+J_m(v)=\mu+J(v)-I_g(v).
\end{equation}
Here the norm of a two-form is the exterior-algebra norm.  Completing a
square gives
\begin{equation}
 \label{prep:electromagnetic-square}
 I_g(v)
 =c_n^2\bigl(|D_m+v^\flat\wedge E^\flat|_g^2
             +|E|_g^2-|v\wedge E|_g^2\bigr)
 \geq c_n^2\bigl(|E|_g^2-|v\wedge E|_g^2\bigr).
\end{equation}
In the purely electric case this is exactly the data convention in
Definition~\ref{def:normalization}.  The auxiliary class will be used
only for already future-timelike data with strict trapping and a strict
charged margin.  No magnetic endpoint theorem is being assumed.

\begin{lemma}[Conformal constraint and expansion formulas]
\label{prep:conformal}
Let $s$ be smooth and set $g_s=e^{2s}g$, $K_s=e^sK$.  If $v$ is a
$g$-unit-ball vector, let $v_s=e^{-s}v$.  Then
\begin{align}
 \label{prep:conformal-constraint}
 e^{2s}\bigl(\mu_s+J_s(v_s)\bigr)
 &=\mu+J(v)-k\Delta_gs-\frac{k(n-2)}2|\dd s|_g^2
       +kK(\nabla s,v),\\
 \label{prep:conformal-expansion}
 e^s\bigl(H_s+\epsilon_a\tr_{S_a,g_s}K_s\bigr)
 &=H+\epsilon_a\tr_{S_a,g}K+k\partial_\nu s.
\end{align}
The second formula holds for either fixed sign $\epsilon_a$.
With $\alpha$ fixed, $E_s=e^{-ns}E$, so in the purely electric case
\begin{equation}
 \label{prep:electric-conformal-cost}
 e^{2s}I_{g_s}(v_s)=c_n^2e^{-2(n-2)s}|E|_g^2.
\end{equation}
In particular, this cost does not increase if $s\geq0$.
\end{lemma}

\begin{proof}
The scalar-curvature and hypersurface conformal laws give
\[
 e^{2s}R_{g_s}=R_g-2k\Delta_gs-k(n-2)|\dd s|_g^2,
 \qquad e^sH_s=H+k\partial_\nu s.
\]
Also $\tr_{g_s}K_s=e^{-s}\tr_gK$.  If
$\Pi=K-(\tr_gK)g$, then $\Pi_s=e^s\Pi$.
Contracting the connection difference
\[
 (\nabla^s-\nabla)^l{}_{ij}
 =s_i\delta_j^l+s_j\delta_i^l-s^lg_{ij}
\]
gives
$J_s=e^{-s}(J+kK(\nabla s,\cdot))$.
These identities prove
\eqref{prep:conformal-constraint}--\eqref{prep:conformal-expansion}.
Finally, $\vol_{g_s}=e^{ns}\vol_g$ and
$i_{E_s}\vol_{g_s}=\alpha$, proving the assertion about $E_s$ and its
norm.
\end{proof}

For $s=2\log U/(n-2)$, the new bulk term in
\eqref{prep:conformal-constraint} is exactly
\begin{equation}
 \label{prep:U-constraint}
 \frac{2k}{(n-2)U}
 \bigl(-\Delta_gU+K(\nabla U,v)\bigr).
\end{equation}
This form of the identity will be used for the end repairs.

\begin{lemma}[Increasing the ADM energy]
\label{prep:end-increase}
For purely electric data with the weak asymptotic hypotheses, choose
$0<\beta<\min\{1,q\}$.  There is a smooth positive function $\phi_0$,
constant on a compact set containing $S$, such that
\[
 \phi_0=r^{2-n}(1-r^{-\beta})\quad\hbox{far out},
 \qquad -\Delta_g\phi_0\geq|K|_g|\dd\phi_0|_g.
\]
For each $t\geq0$, the change
\[
 U_t=1+t\phi_0,\qquad
 g_t=U_t^{4/(n-2)}g,\qquad K_t=U_t^{2/(n-2)}K
\]
with fixed $\alpha$ preserves all numerical hypotheses, all prescribed
trapping signs, and the signed charge.  It satisfies
\begin{equation}
 \label{prep:end-increase-charges}
 \mathcal E_{\mathrm{ADM}}(g_t,K_t)
 =\mathcal E_{\mathrm{ADM}}(g,K)+2t,
 \qquad \mathbf P_{\mathrm{ADM}}(g_t,K_t)
 =\mathbf P_{\mathrm{ADM}}(g,K),
 \qquad a(g_t)\geq a(g).
\end{equation}
If the input has stronger asymptotic derivative bounds, those same
derivative bounds are preserved.  In particular, the lemma applies
with $E=0$ to the strong-tail neutral data used in
Theorem~\ref{thm:neutral-input}.
\end{lemma}

\begin{proof}
For $f(r)=r^{2-n}(1-r^{-\beta})$, asymptotic differentiation gives
\[
 -\Delta_g f
 =\beta(n-2+\beta)r^{-n-\beta}+O(r^{-n-q}),
 \qquad |K|_g|\dd f|_g=O(r^{-n-q}).
\]
Thus the desired strict inequality holds on a sufficiently distant
coordinate region.  On that region $-f'>0$.  Multiply $-f'$ by a smooth
nondecreasing cutoff that changes from zero to one, entirely in this
valid region, and integrate from infinity.  If $\chi$ is the cutoff,
the resulting radial function satisfies
\[
 -\Delta_g\phi_0-|K|_g|\dd\phi_0|_g
 =\chi\bigl(-\Delta_g f-|K|_g|\dd f|_g\bigr)
   +\chi'(-f')|\dd r|_g^2\geq0.
\]
It is positive and constant on the remaining core.  Equations
\eqref{prep:U-constraint} and \eqref{prep:electric-conformal-cost}
preserve charged DEC; the expansion is only rescaled on $S$.
The asymptotic constraints remain integrable, since the new Laplacian
term is $O(r^{-n-\beta})$, the drift term is $O(r^{-n-q})$, and the
gradient-square term is integrable.  All other end bounds persist.

The leading change of the metric is
$4t r^{2-n}\delta/(n-2)$, which has ADM energy $2t$ in the stated
normalization.  The change in the momentum integrand has sphere integral
$O(r^{-q})$, and hence zero limit.  Charge is unchanged because its
flux form is unchanged.  Finally $U_t\geq1$, so every full-cut area
increases.
\end{proof}

\begin{remark}[Order of the endpoint argument]
\label{prep:endpoint-order}
Suppose the numerical estimate has first been proved for all
future-timelike purely electric data, without assuming the conclusion
for other data.  If an original data set had
$\mathcal E_{\mathrm{ADM}}\leq|\mathbf P_{\mathrm{ADM}}|$, apply that
estimate to Lemma~\ref{prep:end-increase} with
\[
 t>t_*:=\frac{|\mathbf P_{\mathrm{ADM}}|
                    -\mathcal E_{\mathrm{ADM}}}{2}.
\]
The invariant mass tends to zero as $t\downarrow t_*$, whereas the
neutral part of the numerical estimate gives the fixed positive lower
bound $X_A/2$.  This contradiction excludes the endpoint.  Thus the
following preparation needs to treat future-timelike data only; the
remark is a conditional reduction, not an invocation of the neutral
input on null data.
\end{remark}

\subsection{Compact inverses for the linear cutoff errors}

\begin{lemma}[Compact inverses with the required moments]
\label{prep:compact-inverses}
Fix nested compact subannuli inside $\{1<|z|<4\}$, with sufficiently
many intermediate subannuli to allow two support enlargements, and fix
$1<p<\infty$.  For smooth sources supported in the innermost subannulus
there are linear constructions with the following properties.
\begin{enumerate}[label=(\roman*)]
\item If $f$ has zero constant and affine moments, there is a compactly
supported symmetric tensor $H$ in the outer subannulus satisfying
\[
 \partial_i\partial_jH_{ij}=f,
 \qquad \|H\|_{W^{3,p}}\leq C\|f\|_{W^{1,p}}.
\]
\item If the vector source $F$ has zero moments against all Euclidean
translations and rotations, there is a compactly supported symmetric
tensor $B$ in the outer subannulus satisfying
\[
 \partial_jB_{ij}=F_i,
 \qquad \|B\|_{W^{2,p}}\leq C\|F\|_{W^{1,p}}.
\]
\end{enumerate}
All supports and constants are fixed independently of the sources.
There are fixed smooth moment projections $\Pi_{\rm aff}$ and
$\Pi_{\rm Kill}$ onto compactly supported bumps such that the preceding
inverses apply to $f-\Pi_{\rm aff}f$ and $F-\Pi_{\rm Kill}F$.
The norms of these projections are bounded by the sums of the absolute
values of the corresponding moments.  In particular, when $p>n$, the
two corrections have controlled $C^{2,\gamma}$ and $C^{1,\gamma}$ norms,
where $\gamma=1-n/p>0$.
\end{lemma}

\begin{proof}
We first describe a compact scalar divergence inverse.  On a ball a
regularized Bogovskii operator, with a fixed
$\rho\in C_c^\infty$ of integral one, is given by
\[
 (\mathcal B_\rho h)(x)
 =\int h(y)(x-y)\int_1^\infty
       \rho(y+t(x-y))t^{n-1}\,\dd t\,\dd y.
\]
It has divergence $h-\rho\int h$ and gains one Sobolev derivative;
see \cite{CostabelMcIntosh2010}.  Its support is contained in the
convex hull of the supports of $h$ and $\rho$.  The derivative bounds
needed here also follow directly from the first-derivative estimate
and
\[
 \partial_j\mathcal B_\rho h
 =\mathcal B_\rho(\partial_jh)
       +\mathcal B_{\partial_j\rho}h,
\]
first for compactly supported smooth $h$.  The kernel estimate does not
require the new kernel $\partial_j\rho$ to have integral one.

To work in a connected annulus, cover the prescribed source support by
finitely many balls compactly contained in a larger subannulus, adding
balls along paths so that the intersection graph is connected.  Choose
a partition of unity and a spanning tree of that graph.  In each
parent--child overlap fix a unit-integral smooth bump.  Transfer the
integral of a child source to its parent using that bump, beginning at
the leaves.  This writes a source of total integral zero as a sum of
mean-zero sources on the balls.  The transfer coefficients are bounded
by the original source norm.  Summing the ball operators gives an
operator $\mathcal B$ with
\begin{equation}
 \label{prep:bogovskii-bounds}
 \operatorname{div}\mathcal B h=h,
 \qquad \|\mathcal B h\|_{W^{j+1,p}}
       \leq C_j\|h\|_{W^{j,p}},\qquad j=0,1,2,
\end{equation}
and support in a fixed larger compact subset of the annulus.

For (i), set $u=\mathcal B f$.  Its component means vanish, since compact
support and integration by parts give
\[
 \int u_i=-\int z_i\partial_j u_j=-\int z_i f=0.
\]
Apply $\mathcal B$ to each $u_i$ to obtain $T_{ij}$ with
$\partial_jT_{ij}=u_i$.  The symmetric part
$H=(T+T^{\mathsf T})/2$ has the same double divergence as $T$.
Two applications of \eqref{prep:bogovskii-bounds} give its $W^{3,p}$
bound.

For (ii), translation moments first give a row primitive $B^0$ with
$\partial_jB^0_{ij}=F_i$ and a $W^{2,p}$ bound.  Write
$A_{ij}=(B^0_{ij}-B^0_{ji})/2$.  Rotation moments imply
\[
 2\int A_{ij}=-\int(z_jF_i-z_iF_j)=0.
\]
Solve $\partial_lD_{ijl}=-A_{ij}$ using \eqref{prep:bogovskii-bounds},
choosing $D_{ijl}=-D_{jil}$.  Thus $D$ has a $W^{3,p}$ bound.  Define
\[
 C_{ij}=\partial_l(D_{ijl}-D_{ilj}-D_{jli}).
\]
Then
\[
 C_{ij}-C_{ji}=-2A_{ij},\qquad
 \partial_jC_{ij}=0.
\]
For the divergence identity the first two terms cancel after exchanging
$j,l$, and the third vanishes by antisymmetry in those indices.
Consequently $B=B^0+C$ is symmetric, has the required divergence, and
has a $W^{2,p}$ bound.  None of these differentiations enlarges support.

For completeness, the moment bumps can be fixed without a choice
depending on the source.  Let $\eta$ be a smooth nonnegative function
supported in the innermost subannulus and positive on a ball.  For the
affine basis $\varphi_a\in\{1,z_1,\ldots,z_n\}$, the matrix
$G_{ab}=\int\eta\varphi_a\varphi_b$ is positive definite.  The functions
$\eta\sum_b(G^{-1})_{ab}\varphi_b$ are dual moment bumps.  For vector
sources use instead a basis of Euclidean Killing fields and its Gram
matrix for the $\eta$-weighted Euclidean inner product.  A nonzero
affine function or Killing field cannot vanish on a ball, so both Gram
matrices are invertible.  These projections have all the asserted norm
and support bounds.  Sobolev embedding finishes the proof.
\end{proof}

\subsection{A strict conformal direction with weak asymptotic data}

\begin{lemma}[The decaying Neumann problem and strictification]
\label{prep:strictification}
Assume $g-\delta=O_2(r^{-q})$ with $q>(n-2)/2$, and let
$0<\beta<\min\{1,q\}$.  Let $B\geq0$ be smooth, with
$B\geq|K|_g$ and $B=O_1(r^{-1-q})$.  It is permissible instead that
$B$ be compactly supported.  There is a unique smooth positive decaying
solution
\begin{equation}
 \label{prep:neumann-equation}
 -\Delta_g\phi
 =B\sqrt{|\dd\phi|_g^2+r^{-2k}}+r^{-n-\beta},
 \qquad \partial_\nu\phi=-1\ \text{on }S,
 \qquad \phi\longrightarrow0\ \text{at infinity}.
\end{equation}
It satisfies
\begin{equation}
 \label{prep:neumann-asymptotics}
 \phi=O_2(r^{2-n}),\qquad
 \frac{-\Delta_g\phi}{r^{-n-\beta}}\longrightarrow1,
 \qquad
 \lim_{R\to\infty}\int_{S_R}\partial_{\nu_R}\phi\,\dd A_g
 \text{ exists and is strictly negative}.
\end{equation}
If $B$ is compactly supported and the metric end has all-order
$r^{2-n}$ bounds, then $\phi=O_l(r^{2-n})$ for every $l$.

For purely electric data satisfying charged DEC and the specified weak
trapping inequalities, the change
\[
 g_\delta=e^{2\delta\phi}g,
 \qquad K_\delta=e^{\delta\phi}K,
 \qquad \alpha_\delta=\alpha
\]
has, for every sufficiently small $\delta>0$, a charged margin
$\mu_{m,\delta}-|J_{m,\delta}|\geq c\delta r^{-n-\beta}$ and strict
trapping for every prescribed sign.  It changes the ADM energy by
$O(\delta)$, leaves momentum and charge unchanged, and changes the
metric by a uniform factor tending to one.  It does not decrease any
full-cut area.  The same conclusion holds for a fixed compactly
supported $D_m$ if there is already a positive charged margin on its
support.
\end{lemma}

\begin{proof}
Write $a_0=r^{-k}$ and $\rho_0=r^{-n-\beta}$.  Truncate at a distant
coordinate sphere $S_R$, and continue in $\lambda\in[0,1]$ the problem
\[
 -\Delta_g\phi_R
 =\lambda B\sqrt{|\dd\phi_R|_g^2+a_0^2}+\rho_0,
 \qquad \partial_\nu\phi_R=-1\text{ on }S,
 \qquad \phi_R=0\text{ on }S_R.
\]
The inner Neumann and outer Dirichlet faces are disjoint.  The Poisson
problem at $\lambda=0$ has a unique smooth solution.  The linearization
at any solution is
\[
 -\Delta_g h-\lambda B
 \frac{\langle\dd\phi_R,\dd h\rangle_g}
      {\sqrt{|\dd\phi_R|_g^2+a_0^2}},
 \qquad \partial_\nu h=0,\quad h|_{S_R}=0.
\]
Its drift is bounded by $B$; the maximum principle and boundary Hopf
lemma give a trivial kernel.  Fredholm theory for this scalar mixed
boundary problem, or continuation from the mixed Laplacian, then gives
invertibility.  We use the usual scalar local and boundary estimates
in this argument; see \cite{GilbargTrudinger2001}.  There is no meeting
edge between Dirichlet and Neumann conditions.

Every solution is nonnegative.  An interior negative minimum contradicts
the positive right side, and a minimum on $S$ contradicts the inward
derivative $-1$.  Given a supremum bound, local and boundary
$W^{2,p}$ estimates, interpolation, and the at-most-linear gradient
growth of the equation give uniform $W^{2,p}$ bounds on a fixed
truncation.  More explicitly, for $p>n$,
\[
 \|\phi_R\|_{W^{2,p}}
 \leq C_R(1+\|\dd\phi_R\|_{L^p}+\|\phi_R\|_{L^p})
 \leq\tfrac12\|\phi_R\|_{W^{2,p}}
       +C_R'(1+\|\phi_R\|_\infty).
\]
The resulting $C^{1,\gamma}$ bound makes the right side
$C^{0,\gamma}$, and Schauder estimates and differentiation give
smooth bounds.  Positivity of $a_0$ makes the square root a smooth
function of the gradient on every compact set.

We next obtain the required supremum bounds.  A normalized blow-up on
a fixed truncation, if those bounds failed, would give a nonzero
nonnegative solution $v$ of
\[
 -\Delta_g v=\lambda B|\dd v|,
 \qquad \partial_\nu v=0,
 \qquad v|_{S_R}=0.
\]
The preceding estimates apply to the normalized solutions as well.
The limiting equation is homogeneous with bounded drift: where
$\dd v\ne0$ use $\lambda B\nabla v/|\dd v|$, and set the drift to
zero elsewhere.  The maximum principle and Hopf lemma exclude such a
$v$.  This closes the continuation on every fixed truncation.

For uniformity as $R\to\infty$, use
$f_0(r)=r^{2-n}(1-r^{-\beta})$ on a sufficiently distant fixed region.
As in Lemma~\ref{prep:end-increase},
\[
 -\Delta_g f_0-B|\dd f_0|\geq c_0\rho_0,
 \qquad Ba_0=O(r^{-n-q})\leq C_0\rho_0.
\]
For $A$ larger than a fixed constant and chosen also to dominate the
values on a fixed inner sphere $S_{R_0}$, $Af_0$ is a supersolution,
because
\[
 \lambda B\sqrt{A^2|\dd f_0|^2+a_0^2}+\rho_0
 \leq AB|\dd f_0|+Ba_0+\rho_0
 \leq-\Delta_g(Af_0).
\]
Subtraction of the two equations produces a bounded-drift inequality,
so comparison is valid without a gradient monotonicity assumption.
The outer Dirichlet inequality is automatic.  It follows that
\begin{equation}
 \label{prep:neumann-tail-comparison}
 0\leq\phi_R(x)
 \leq C\bigl(1+\sup_{S_{R_0}}\phi_R\bigr)r^{2-n},
 \qquad R_0\leq r\leq R,
\end{equation}
with $C$ independent of $R$ and $\lambda$.

If the suprema on an exhaustion were unbounded, divide by those
suprema.  Estimate~\eqref{prep:neumann-tail-comparison} forces their
unit maxima into a fixed compact set and bounds their tails by a fixed
multiple of $r^{2-n}$.  Local estimates, including the fixed inner
Neumann charts, yield a nonzero decaying limit solving the same
homogeneous bounded-drift equation with zero Neumann data on $S$.
Its positive maximum is attained either in the interior or on $S$;
the strong maximum principle and Hopf lemma exclude both alternatives.
This proves a uniform supremum bound.  Exhaustion and local compactness
now produce a smooth solution of \eqref{prep:neumann-equation}.
It is strictly positive by the same minimum argument.  The difference
of two decaying solutions solves a homogeneous equation with bounded
drift and zero Neumann data, giving uniqueness.

To check the asymptotic derivatives using only the weak hypotheses,
rescale on fixed end annuli by
$\phi_R^*(z)=R^{n-2}\phi(Rz)$ and $g_R(z)=g(Rz)$.
The equation becomes
\[
 -\Delta_{g_R}\phi_R^*
 =RB(Rz)\sqrt{|\dd\phi_R^*|_{g_R}^2+|z|^{-2k}}
       +R^{-\beta}|z|^{-n-\beta}.
\]
Its solutions have uniformly bounded suprema by
\eqref{prep:neumann-tail-comparison}.  The two controlled derivatives
of $g$ give uniformly Lipschitz first-order coefficients and uniformly
H\"older principal coefficients on these rescaled annuli.
$W^{2,p}$ estimates followed by $C^{1,\gamma}$ and Schauder estimates
therefore give uniform second-derivative bounds.  This proves
$\phi=O_2(r^{2-n})$ without a bound on a third metric derivative.
If all end derivatives of $g$ are controlled and $B=0$ there, the same
argument differentiates repeatedly.  The equation and these bounds
also give
\[
 B\sqrt{|\dd\phi|^2+r^{-2k}}=O(r^{-n-q}),
\]
which proves the ratio limit in
\eqref{prep:neumann-asymptotics}.  The source is integrable.  The
divergence theorem, remembering that the outward domain normal on
$S$ is $-\nu$, gives
\begin{equation}
 \label{prep:neumann-flux}
 \int_{S_R}\partial_{\nu_R}\phi\,\dd A_g
 =-\Area_g(S)-\int_{\Omega\cap\{r\leq R\}}
       \bigl(B\sqrt{|\dd\phi|^2+r^{-2k}}+\rho_0\bigr)\,\dd V_g.
\end{equation}
The flux thus has a finite negative limit.  Its difference from the
Euclidean derivative flux is $O(R^{-q})$.

For the final strictification, apply
\eqref{prep:conformal-constraint} with $s=\delta\phi$.
The new linear bulk term obeys
\[
 k\delta\bigl(-\Delta_g\phi+K(\nabla\phi,v)\bigr)
 \geq k\delta\rho_0.
\]
Moreover
\[
 |\dd\phi|^2/\rho_0=O(r^{2-n+\beta})
\]
is globally bounded.  Thus for sufficiently small $\delta$ the
quadratic term costs at most $k\delta\rho_0/2$.  In the purely electric
case the conformal electromagnetic cost decreases.  On a fixed
magnetic support any existing positive margin absorbs the small
change in \eqref{prep:electromagnetic-cost}; outside that support the
purely electric argument applies.  This proves a positive multiple
of $\delta\rho_0$ as a charged margin.  On every original boundary
component, exactly,
\[
 H_\delta+\epsilon_a\tr_{S_a,g_\delta}K_\delta
 =e^{-\delta\phi}
       \bigl(H+\epsilon_a\tr_{S_a,g}K-k\delta\bigr)<0.
\]
The energy change is
\begin{equation}
 \label{prep:strictification-energy}
 \mathcal E_{\mathrm{ADM}}(g_\delta,K_\delta)
       -\mathcal E_{\mathrm{ADM}}(g,K)
 =-\frac\delta\omega
       \lim_{R\to\infty}\int_{S_R}\partial_{\nu_R}\phi\,\dd A_g.
\end{equation}
Terms quadratic in the decaying conformal factor have vanishing ADM
flux.  The momentum change again has zero limit, since its integrand
has one extra factor $\phi=O(r^{2-n})$.  Charge is fixed by $\alpha$.
Finally $\phi$ is positive and bounded, giving both the area
monotonicity and the uniform relative metric convergence.  On a
prepared end $K=B=0$, and the scalar conformal formula with
$-\Delta\phi=\rho_0$ also proves the asserted prepared scalar decay
$O(r^{-n-\beta})$.
\end{proof}

\begin{corollary}[Strict direction for compact variations]
\label{prep:strict-direction}
On the original weakly asymptotically flat exterior, choose a smooth
majorant $B\geq|K|$ equal to $Cr^{-1-q}$ far out and apply
Lemma~\ref{prep:strictification}.  Consider the tangent direction
\[
 (g',K',\alpha',D_m')=(2\phi g,\phi K,0,0)
\]
at purely electric data, transporting $v$ by the metric square root.
On an active ray $\mu_m+J_m(v)=0$, the derivative
$\mathfrak C'=2(\mu_m+J_m(v))'$ satisfies
\begin{equation}
 \label{prep:strict-ray-derivative}
 \mathfrak C'
 =2k\bigl(-\Delta\phi+K(\nabla\phi,v)\bigr)
       +4(n-2)c_n^2\phi|E|^2
 \geq2k\rho_0,
 \qquad \mathfrak C'/\rho_0\longrightarrow2k.
\end{equation}
At a marginal boundary $\theta_+'=-k$.  The energy derivative is finite,
the momentum derivative is zero, and $\phi=O_2(r^{2-n})$.
\end{corollary}

\begin{proof}
A smooth radial majorant on the end can be patched to a sufficiently
large smooth function on the core; it need not reproduce higher
derivatives of $K$.  Differentiate
\eqref{prep:conformal-constraint} and
\eqref{prep:electric-conformal-cost}.  The derivative of the overall
factor $e^{-2\delta\phi}$ multiplies the active value zero.  The drift
term is $O(r^{-n-q})$, and the electric term divided by $\rho_0$ tends
to zero.  The assertions now follow from
\eqref{prep:neumann-asymptotics},
\eqref{prep:conformal-expansion}, and
\eqref{prep:strictification-energy}.
\end{proof}

\subsection{Matching a neutral vacuum end and moving it to rest}

We state the reduction in a form that includes the small magnetic
variations needed later.  The purely electric case requires no strict
margin in the original data.

\begin{proposition}[Preparation at rest]
\label{prep:rest-preparation}
Suppose the original data have a future-timelike ADM vector, and set
$m=(\mathcal E_{\mathrm{ADM}}^2-|\mathbf P_{\mathrm{ADM}}|^2)^{1/2}>0$.
Choose once and for all
\[
 \frac{n-2}{2}<q<n-2,
 \qquad 0<\beta<\min\{1,q\},
\]
weakening the original decay exponent if necessary.  There is a sequence
of smooth data on the same exterior, with the same electric flux form
and charge, satisfying
\begin{align}
 \label{prep:prepared-properties}
 &K_j\text{ is compactly supported},\qquad
   g_j-\delta=O_l(r^{2-n})\quad\text{for every }l,\\
 &R_{g_j}=O(r^{-n-\beta}),\qquad
   \mu_{m,j}-|J_{m,j}|_{g_j}\geq c_jr^{-n-\beta},\qquad
   H_j+\epsilon_a\tr_{S_a,g_j}K_j<0,
 \nonumber
\end{align}
where $c_j>0$.  Its momentum vanishes and
\begin{equation}
 \label{prep:prepared-limits}
 \mathcal E_{\mathrm{ADM}}(g_j,K_j)\longrightarrow m,
 \qquad \liminf_{j\to\infty}a(g_j)\geq a(g).
\end{equation}
The same conclusion holds with a fixed compactly supported closed
$D_m$, provided the original data are already future-timelike, have
strict trapping, and satisfy a strict charged margin
$\mu_m-|J_m|\geq c r^{-n-\beta}$ for some $c>0$.
The constants and support radii in
\eqref{prep:prepared-properties} need not be uniform in $j$.
\end{proposition}

\begin{proof}
The proof has four stages.  The last, a small strictification, is supplied
by Lemma~\ref{prep:strictification}.  We first construct data with
nonnegative charged margin, compactly supported second tensor, and
energy approaching $m$.

\paragraph{The matching vacuum plane.}
In isotropic coordinates $(b,y)$, the Schwarzschild--Tangherlini metric
of mass $m$ has spatial factor
$(1+m|y|^{2-n}/2)^{4/(n-2)}$.  Its leading spacetime perturbation from
Minkowski space is
\[
 2m|y|^{2-n}\left(\dd b^2+\frac{|\dd y|^2}{n-2}\right),
\]
with remainders of squared order and their derivatives.  After rotating
the original end coordinates, choose $v\in[0,1)$ and
$\gamma=(1-v^2)^{-1/2}$ so that the desired ADM vector is
$(\gamma m,\gamma mv,0,\ldots,0)$.  Put
\[
 b=\gamma(T+vx^1),\qquad y^1=\gamma(x^1+vT),\qquad
 y^\perp=x^\perp,
\]
and take the plane $T=0$ with the future normal.  Its induced metric has
leading spatial perturbation
\[
 h^{\mathrm{pl}}_{ij}
 =\frac{2m}{n-2}|y|^{2-n}
   \bigl(\delta_{ij}+k(\gamma^2-1)\delta_{i1}\delta_{j1}\bigr).
\]
Direct differentiation gives the leading energy flux vector
$2km\gamma^2x_i|y|^{-n}$.  With the convention
$K=G(\nabla N,\cdot)$, the linear expression
\[
 2K_{ij}=\partial_Tg_{ij}-\partial_i g_{0j}-\partial_jg_{0i}
\]
gives first momentum row $km\gamma^2v x_i|y|^{-n}$; the transverse
fluxes integrate to zero by reflection.  The identity
\[
 \frac1\omega\int_{S^{n-1}}|A\theta|^{-n}\,\dd A_\theta
 =\frac1{\det A},\qquad A=\operatorname{diag}(\gamma,1,\ldots,1),
\]
follows by computing the volume of the ellipsoid $|Ax|\leq1$ in polar
coordinates.  Hence the plane has exactly the specified energy and
momentum with the problem's ADM normalizations.  This is a neutral
vacuum model; the electric form will subsequently be retained on the
underlying exterior, not induced from a charged boosted slice.

\paragraph{Cutoff errors and their moments.}
Let $L\to\infty$ and work on $1<|z|<4$, $x=Lz$.  Use a fixed cutoff,
constant near both boundary spheres, to interpolate the components
$g(Lz)$ and $LK(Lz)$ to the matching plane.  The flat linear constraint
operators on a metric perturbation $h$ and a tensor $p$ are
\[
 \mathcal L h=\partial_i\partial_jh_{ij}-\Delta\tr_\delta h,
 \qquad (\mathcal Mp)_i=\partial_j(p_{ij}-(\tr_\delta p)\delta_{ij}).
\]
Relative to the cutoff interpolation of the scaled physical sources
$(2\mu,J)$, the new sources are the flat cutoff commutators plus a
$C^0$ remainder $O(L^{-2q})$.  The scalar commutator consists of
$(D\chi)(Dh)$ and $(D^2\chi)h$; its $W^{1,p}$ norm is
$O(L^{-q})$ using only the given $C^2$ bounds on $g$.
The vector commutator is $(D\chi)(p-(\tr p)\delta)$ and has the same
$W^{1,p}$ bound using only the given $C^1$ bounds on $K$.  We do not
differentiate the nonlinear remainder or the physical constraint
sources.

All affine scalar moments and Killing vector moments of these
commutators are $o(L^{2-n})$.  Here is the reason that no center-of-mass
or angular-momentum limit is required.  In physical coordinates the
difference between a flat linear source and its physical source is
$O(r^{-2q-2})$, which is integrable.  The physical sources are
integrable by hypothesis, and the matching vacuum end has integrable
linear sources for the same reason.  For any such integrable source
$F$,
\begin{equation}
 \label{prep:first-moment-oL}
 \int_{r<4L}r|F|\,\dd x=o(L).
\end{equation}
Indeed split at a fixed radius, divide by $L$, and then make the
integrable tail small.  Green's formula for $\mathcal L$ against an
affine function and for $\mathcal M$ against a Killing field therefore
bounds the corresponding linear boundary fluxes by $o(L)$.
For constant scalar and translation tests, the difference of boundary
fluxes tends to zero by the matched ADM vector.  The difference between
the physical and flat momentum boundary integrands has integral
$O(r^{n-2-2q})=o(1)$, so the same matching applies to $\mathcal M$.
Integrating a commutator against a test leaves these boundary terms
minus a cutoff-weighted source integral.  For a degree-one test the
latter, divided by $L$, tends to zero by
\eqref{prep:first-moment-oL}; for a constant test it tends to zero by
integrability.  Rescaling $\dd x=L^n\dd z$ and the second-order source
by $L^2$ proves the claimed $o(L^{2-n})$ moments.

Apply Lemma~\ref{prep:compact-inverses} to cancel the commutators after
their moment projections have been removed.  Invert the full trace
reversal $h\mapsto h-(\tr h)\delta$ and
$p\mapsto p-(\tr p)\delta$ after constructing the two symmetric
primitives.  This changes no norm or support conclusion.  Taking
$p>n$, the resulting corrections have scaled $C^2,C^1$ norms
$O(L^{-q})$.  They are smooth and supported strictly inside the
transition annulus.  The remaining moment bumps are
$o(L^{2-n})$ in these fixed-scale norms.  Nonlinear constraint errors
are still $O(L^{-2q})$; comparison of the original and interpolated
unit balls incurs the same order.  In physical coordinates the charged
DEC deficit on this annulus is consequently
\begin{equation}
 \label{prep:annular-deficit}
 o(L^{-n}).
\end{equation}
The retained electric field contributes at most $O(L^{-2k})$, also
$o(L^{-n})$.  For the auxiliary magnetic class its compact support is
inside the untouched core once $L$ is large.  Positive definiteness
follows from the small scaled metric correction.  Only the two weak
asymptotic derivative bounds were used in this construction.

\paragraph{Bending and a conformal repair.}
Beyond the gluing annulus, use $y$ as spatial parameters on the neutral
vacuum model.  The plane is $b=vy^1$.  Multiply this function by a cutoff
that decreases slowly as a function of $\log(|y|/L)$, and that is zero
after a fixed sufficiently long interval in that variable.  Its
derivative can be chosen so small, in terms of the fixed $v<1$, that
the Euclidean slope stays strictly below one.  For large $L$ the graph
is then spacelike in the exact vacuum model as well.  It joins the
matching plane to the static slice $b=0$ within a fixed scaled annulus.
The resulting induced metric is uniformly Euclidean-comparable there,
with
\[
 |Dg|+|K|\leq C/r.
\]
The constants may depend on the original timelike vector; no uniformity
as that vector tends to null is required.  The graph is exactly vacuum,
so its only charged deficit is the cost of the retained electric form.

We repair both deficits by \eqref{prep:U-constraint}.  Choose numbers
$\delta_L\downarrow0$ dominating the coefficient in
\eqref{prep:annular-deficit}.  Choose next $\varepsilon_L\downarrow0$
so slowly that
\[
 \delta_L/\varepsilon_L\longrightarrow0,
 \qquad L^{2-n}/\varepsilon_L\longrightarrow0.
\]
On the fixed scaled radial interval containing the transitions, let
$s=|y|/L$ and construct a nonnegative smooth function $z_L$ that starts
at zero and satisfies
\[
 z_L'=C_0z_L+\chi_*,
\]
where $\chi_*$ is a fixed nonnegative bump equal to one over the
possible annular deficits.  Start its support strictly before those
deficits.  The constant $C_0$ is chosen using the uniform radial
ellipticity and the bounds on $Dg,K$.  If $F_L'=-z_L$, then on this
interval
\[
 -\Delta_g\bigl(F_L(|y|/L)\bigr)
       -|K|\,|\dd(F_L(|y|/L))|
 \geq cL^{-2}\chi_*.
\]
The same inequality is nonnegative on the initial ramp; it follows
directly by separating the $z_L'$ coefficient, bounded below by radial
ellipticity, from terms bounded by a constant times $z_L$.

On the static part put
\[
 A_L(s)=1+\frac{m}{2(Ls)^{n-2}}.
\]
Continue $z_L$ smoothly so that the radial Laplace flux
$P_L(s)=s^kA_L(s)^2z_L(s)$ is increasing, has derivative bounded below
by a positive multiple of $s^{-k}$ on the continuation, and satisfies
$P_L'(s)=s^{-k}$ sufficiently far out.  This can be done by blending
positive derivatives after the end of the preceding interval.  All
profiles and their fixed-order derivatives remain uniformly bounded
on compact scaled intervals.  Set
\[
 F_L(s)=\int_s^\infty z_L(\sigma)\,\dd\sigma,
 \qquad U_L=1+\varepsilon_LL^{2-n}F_L(|y|/L).
\]
Thus $F_L>0$, is constant before the repair region, and has a tail
$a_Ls^{2-n}+O(s^{2-2k})$ with uniformly bounded $a_L>0$.
On the static end the exact radial Laplacian formula is
\[
 -\Delta_g U_L
 =\varepsilon_LL^{-n}A_L(s)^{-2n/(n-2)}s^{-k}P_L'(s)
 \geq c\varepsilon_LL^{-n}s^{-2k}.
\]
It pays for the retained electric cost $C(Ls)^{-2k}$ by the second
condition on $\varepsilon_L$.  The repair on the transition pays for
\eqref{prep:annular-deficit} by the first condition.  Uniform
comparability of $U_L$ to one allows the fixed constants in
\eqref{prep:U-constraint} to be included in these choices.

In the untouched core $U_L$ is constant.  Purely electric DEC is
preserved there for every $U_L\geq1$.  In the auxiliary magnetic class,
the original positive margin on the fixed magnetic support absorbs
the $o(1)$ change of the electromagnetic cost.  Outside that support
the purely electric argument applies.  The trapping inequalities are
preserved on the original boundary, since there $U_L$ is constant.
The repaired data now have compactly supported $K$ and an all-order
static conformal end.  Their final ADM energy, computed in the static
$y$ coordinates, is
\[
 m+2\varepsilon_La_L=m+o(1),
\]
and their momentum is zero.  The scalar curvature on the static end is
$O(r^{-2k})$.  The electric form remains exactly $\alpha$ throughout.

\paragraph{Comparison of the enclosing areas.}
Here no minimizing hypersurface is selected.  In the original $x$
coordinates the repaired metric has a fixed positive Euclidean lower
bound throughout the modified end.  Choose a sufficiently small fixed
$c>0$ and a radial diffeomorphism $\Psi_L$ that is the identity for
$r\leq\sqrt L$, whose radial derivative decreases smoothly to $c$
by $r=2\sqrt L$, and whose radial function is $cr+O(\sqrt L)$
thereafter.  It is the identity on the compact part and preserves
orientation.  Up to radius $L$, the original and repaired metrics agree
apart from the nondecreasing conformal factor; the asymptotic
Euclidean errors on $r\geq\sqrt L$ are $o(1)$, and both radial and
tangential derivatives of $\Psi_L$ are at most one.  On $r\geq L$
those derivatives are at most $c+O(L^{-1/2})$.  Choosing $c$ using the
fixed lower bound of the repaired metric therefore gives the global
quadratic-form comparison
\begin{equation}
 \label{prep:area-metric-comparison}
 \Psi_L^*g\leq(1+o(1))g_L^{\mathrm{rep}}.
\end{equation}
The pullback evaluates $g$ only at radii tending to infinity in this
last region, so its own Euclidean comparison error also tends to zero.
The map $\Psi_L$ carries full cuts bijectively to full cuts, fixing the
original boundary.  Restricting
\eqref{prep:area-metric-comparison} to each tangent $k$-plane and taking
the infimum gives
\[
 a(g_L^{\mathrm{rep}})\geq(1-o(1))a(g).
\]
This argument applies to disconnected cuts and coincident boundary
portions, and is independent of any regularity or attainment of the
infimum.

Finally apply Lemma~\ref{prep:strictification} with a compact drift
majorizing the now compactly supported tensor.  For each fixed $L$,
take the strictification parameter positive and sufficiently small
that the additional energy change and the uniform relative metric
change are less than $L^{-1}$.  This order of choices is allowed even
if the constants of that lemma depend on $L$.  A diagonal sequence
gives \eqref{prep:prepared-properties}--\eqref{prep:prepared-limits}.
The strictification keeps $K$ compactly supported, preserves $\alpha$,
and makes all prescribed signs strict simultaneously.  Since $m>0$,
the final energies are positive for all sufficiently large members of
the sequence.
\end{proof}

\subsection{Flux forms on a later filling and the final reduction}

\begin{lemma}[Extension with sources confined to the filling]
\label{prep:source-extension}
Suppose a compact filling and inner collars are later attached to the
original boundary, with no constraint imposed behind that boundary.
The smooth electric flux form $\alpha$ extends smoothly across the
boundary, agrees with the original form throughout $\Omega$, and can
be made zero beyond a fixed inner collar.  Its exterior derivative is
supported strictly on the filling side.  In a family obtained by
lengthening product necks beyond fixed collars, the extension has
uniform fixed-order bounds on those collars and is zero on the
additional long portions.
\end{lemma}

\begin{proof}
In collar coordinates with $s\geq0$ on the original side, write
$\alpha=a(s)+\dd s\wedge b(s)$, where $a(s)$ and $b(s)$ are tangential
forms.  Closedness says $\dd_Sa(s)=0$ and
$\partial_sa(s)=\dd_Sb(s)$.  Extend the smooth family $b(s)$ across
zero and define on the interior collar
\[
 a(s)=a(0)+\int_0^s\dd_Sb(t)\,\dd t.
\]
This gives a closed smooth extension across $S$, with exactly the
original jets.  Multiply it by a cutoff on the interior side that is
one on a smaller collar and zero before the end of the fixed collar,
then extend by zero farther into the filling.  The derivative of the
result is supported where that cutoff changes, a positive collar
distance behind $S$.  The same fixed extension and cutoff work before
every added neck, proving uniformity.  A globally closed extension is
neither asserted nor needed; for nonzero total charge it would in
general contradict Stokes on a one-ended filling.
\end{proof}

\begin{proposition}[Reduction of the numerical theorem]
\label{prep:prepared-reduction}
Assume that for every prepared data set in
\eqref{prep:prepared-properties}, with positive energy and zero
momentum, the construction of Proposition~\ref{flux:prepared-bound}
establishes, for every $0<s<1$,
\begin{equation}
 \label{prep:prepared-bound-hypothesis}
 \mathcal E_{\mathrm{ADM}}
 \geq s|Q|+\frac{1-s^2}{2}
                    \left(\frac{a}{\omega}\right)^{(k-1)/k}.
\end{equation}
Then Theorem~\ref{thm:numerical} follows for the original purely electric
data.  The same inequality with invariant mass in place of energy
holds on the strict, future-timelike auxiliary magnetic class described
above.
\end{proposition}

\begin{proof}
First fix a future-timelike original data set and a number $s\in(0,1)$.
Apply Proposition~\ref{prep:rest-preparation}.  The prepared energies
converge to its invariant mass and are eventually positive; charges
are identical and the full-cut infima have the stated lower limit.
Passing to the limit in \eqref{prep:prepared-bound-hypothesis} gives
\begin{equation}
 \label{prep:invariant-one-parameter}
 m\geq s|Q|+\frac{1-s^2}{2}X_A,
 \qquad 0<s<1.
\end{equation}
Only after this timelike result has been established do we apply
Remark~\ref{prep:endpoint-order}; Lemma~\ref{prep:positive-area} supplies
its strictly positive lower bound.  It proves
$\mathcal E_{\mathrm{ADM}}>|\mathbf P_{\mathrm{ADM}}|$ for every
original purely electric data set, including the initially possible
past or null cases.  Thus \eqref{prep:invariant-one-parameter} applies
to all of them.

Letting $s\uparrow1$ gives $m\geq|Q|$.  If $X_A>|Q|$, the maximum of
the right side of \eqref{prep:invariant-one-parameter} is attained at
$s=|Q|/X_A$ when $Q\ne0$, and by $s\downarrow0$ when $Q=0$; its value
is $(X_A+Q^2/X_A)/2$.  If $X_A\leq|Q|$, the endpoint $s\uparrow1$
gives precisely the required branch.  These alternatives are
equivalent to
$X_A\leq m+\sqrt{m^2-Q^2}$ together with $m\geq|Q|$.
The argument does not assert the polynomial inequality on the
unprescribed branch.  For the auxiliary magnetic class the preparation
already assumes strictness and timelikeness; the first limiting
argument proves its required extension without making a magnetic
endpoint claim.
\end{proof}
